\documentclass[11pt]{article}

\usepackage[margin=1in]{geometry}
\usepackage[T1]{fontenc}
\usepackage[utf8]{inputenc}
\usepackage{lmodern}
\usepackage{microtype}
\usepackage{amsmath,amssymb,amsthm,mathtools}
\usepackage{enumitem}
\usepackage{booktabs}
\usepackage{algorithm}
\usepackage{algpseudocode}
\usepackage{xurl}
\usepackage[hidelinks]{hyperref}

\newtheorem{theorem}{Theorem}
\newtheorem{definition}{Definition}
\newtheorem{remark}{Remark}

\newcommand{\E}{\mathbb{E}}

\title{Certified Menus for Anonymous DHT Lookups\\
\large Distribution-free parameter selection with simultaneous guarantees under many-testing}
\author{}
\date{March 22, 2026 (archive v565)}

\begin{document}
\maketitle

\begin{abstract}
Anonymous DHT lookups expose a familiar tension: stronger mechanisms (longer delegation walks, dummy queries, padding, retries) raise latency and bandwidth, so deployments tune parameters against traces.
Unfortunately, tuning over many candidate configurations is a form of \emph{many-testing}: after the fact, a surveillance adversary (or an unlucky deployment) can ``pick the winner'' among failed configurations, invalidating naive privacy claims.

We give a practical certify-and-select method with a distribution-free finite-sample guarantee that is \emph{simultaneous} over an entire finite policy class.
The output is a \emph{certified menu}: a set of policies that provably satisfy anonymity, availability, and resource SLAs, together with a cost--privacy frontier that operators can switch between without re-running large experiments.
We also show how to certify such menus off-policy from randomized shadow-routing logs and how to decouple expensive privacy labels from cheap telemetry via a two-sample protocol.
The guarantee is deliberately conditional: it certifies the declared witness/interface and logging experiment, not every possible traffic-analysis, padding, token, or deployment environment.
This is an interface note and is intentionally compact: the raw-$M$ menu theorem already yields a complete public claim for a fixed witness/interface, Certified~B only tightens that accounting by replacing $M$ with an effective multiplicity $M_{\mathrm{eff}}$, and Certified~C develops the deployment pipeline (anytime-valid monitoring, slice discovery, and change-control triggers) that consumes certified menus as inputs.
\end{abstract}

\paragraph{Scope.}
We isolate and formalize \emph{Certified Menus for Anonymous DHT Lookups} as a reusable contract surface in the anonymity / Anonymous-DHT series.

\paragraph{Imports.}
This paper owns the many-testing-safe \emph{selection} wrapper.
We treat (i) what privacy labels mean as imported (Anonymity~A; Synthesis~5),\cite{series:oddsinflation,series:endpointbridge}
(ii) state-dependent pitfalls as imported (AnonDHT State~1),\cite{series:state1}
and (iii) notarization mechanics as imported (Evaluation series).
Menus are exported as contract objects (Synthesis~0) and should be referenced by ABOM/OINL manifests (Anonymity~C).\cite{series:contractobjects,series:abomoinl}

\paragraph{Exports.}
A \emph{certified menu} object (family of transcript projections + simultaneous privacy bounds), together with the raw-pool multiplicity accounting needed to publish a complete menu-level guarantee even before any optional signature compression is applied.

\paragraph{Boundary.}
This paper stops once a finite policy class, a shared witness interface, and simultaneous upper bounds have produced a menu object.
Exact public/on-request packet wiring for one exported certified-menu card belongs to the maintained worked example, and theorem-root-to-review routing across that already-fixed card belongs to the maintained series-spine bridge.
Effective-multiplicity compression belongs to Certified~B, disagreement-gated drift monitoring and change-control belong to Certified~C, evaluator notarization belongs to Eval~1, and public lineage / release packaging belong to ABOM/OINL plus Release~A.\cite{series:workedexample,series:seriesspines,series:sigcomp,series:recert,series:eval1,series:releaseA}

\paragraph{Later-terse-paper citation posture.}
When a later short paper only needs to answer \emph{which simultaneous menu contract was declared for this fixed witness/interface?}, it should stop at the minimal public certified-menu card below plus the tiny raw-$M$ arithmetic drill.
It should widen only when multiplicity compression (Certified~B), disagreement-triggered recertification (Certified~C), evaluator notarization (Eval~1), or release-lineage packaging (ABOM/OINL and Release~A) becomes live.
If the live issue is instead the exact public/on-request packet carrying one exported certified-menu card, widen only to the maintained worked example; if it is theorem-root-to-review routing across that already-fixed card, widen only to the maintained series-spine bridge.\cite{series:sigcomp,series:recert,series:eval1,series:abomoinl,series:releaseA,series:workedexample,series:seriesspines}

\paragraph{Artifact policy.}
This archive is source-only; supporting artifacts are available on request. See \cite{series:artifactpolicy}.

\section{Problem statement}
Let $\Pi=\{\pi_1,\dots,\pi_M\}$ be a finite set of lookup policies.
Each policy induces a distribution over transcripts and therefore over a chosen privacy label $L(\pi)$ (e.g., an odds-inflation bound for a projection) and over performance metrics such as latency $T(\pi)$ and bandwidth $B(\pi)$.
We want to publish a \emph{switchable} set of policies that satisfy:
\[
\textsf{Privacy}(\pi)\le \tau_P,\qquad \textsf{Latency}(\pi)\le \tau_T,\qquad \textsf{Bandwidth}(\pi)\le \tau_B,
\]
\emph{simultaneously} for all policies we certify, with a single overall failure probability $\alpha$.

Because the policy pool can be large, post-hoc tuning is a form of \emph{multiple testing}: a failed configuration can be ``found'' after the fact unless the certificate controls error uniformly over the whole pool.
We therefore design the certificate as a family-wise (or simultaneous) guarantee using standard multiplicity control (e.g., Bonferroni/Holm-style corrections) rather than a single-policy bound.\cite{holm1979}

The theorem below does not by itself define the adversary.
The privacy label, transcript projection, logging rule, and witness/interface version must already encode the observation surface being protected.
A menu certificate for committee-contact-count buckets is therefore not a blanket claim about end-to-end timing correlation, browser fingerprinting, padding-machine behavior, token issuance/redemption linkage, or any other unpinned deployment channel.
Those channels either have to be included in the label/interface or explicitly named as outside the certified claim.

\section{Certified menus}
\begin{definition}[Certified menu]
A certified menu at level $\alpha$ is a subset $\widehat\Pi\subseteq\Pi$ produced from data such that with probability at least $1-\alpha$, every $\pi\in\widehat\Pi$ satisfies the published SLAs.
\end{definition}

The key is that the guarantee is \emph{uniform over the output set}.
This blocks the ``pick the winner'' failure mode.

\subsection{A simple, distribution-free construction}
Assume for each policy $\pi_m$ we can compute a per-policy valid upper confidence bound (UCB) for each SLA quantity from $n$ i.i.d. trials:
\[
\Pr\big[\textsf{Privacy}(\pi_m)\le \mathrm{UCB}_P(m)\big]\ge 1-\alpha_P(m)
\]
(and similarly for latency/bandwidth).
These can be simple: e.g., for bounded witnesses use Hoeffding; for Bernoulli ``failure'' events use Clopper--Pearson.

\begin{theorem}[Union-bound menu validity]
If $\sum_{m=1}^M (\alpha_P(m)+\alpha_T(m)+\alpha_B(m))\le \alpha$ and we define
\[
\widehat\Pi := \{\pi_m:\ \mathrm{UCB}_P(m)\le \tau_P,\ \mathrm{UCB}_T(m)\le \tau_T,\ \mathrm{UCB}_B(m)\le \tau_B\},
\]
then $\widehat\Pi$ is a certified menu at level $\alpha$.
\end{theorem}

\begin{remark}[Effective multiplicity]
This theorem is intentionally crude but already complete as a public baseline: publishing the raw policy count $M$, the $\alpha$-split, and the witness-version pin is enough to make the menu claim auditable. Certified~B is an optional tightening that replaces $M$ by an effective multiplicity $M_{\mathrm{eff}}$ when many policies share the same observable behavior.
\end{remark}

\subsection{Switching without re-certifying}
A practical menu includes a cost frontier:
operators might switch among policies in $\widehat\Pi$ as load varies.
Because the guarantee is uniform over the entire certified set, switching does not invalidate the claim.

\section{Off-policy certification from shadow routing}
Often we have logs from a \emph{randomized} policy $\pi_0$ that ``explores'' a superset of routes/actions.
If for each candidate policy $\pi_m$ we can compute an importance weight $w_m(X)$ for a transcript sample $X$ (with a valid concentration argument), then standard off-policy estimators can produce valid bounds.\cite{huang2021offpolicy,kuzborskij2021confident}
This sentence hides the riskiest deployment assumptions, so a release packet should treat them as hard evidence rows rather than prose:
\begin{enumerate}[leftmargin=*]
\item \textbf{Support / overlap:} every action or route fragment that a certified policy may take has positive logged probability under the shadow-routing policy, and the propensity or replay weight used to compute $w_m$ is retained in the attempt log.
\item \textbf{Bounded or audited weights:} if weights are clipped, self-normalized, or otherwise stabilized, the resulting bias/coverage tradeoff is part of the public certificate rather than silently folded into the UCB.
\item \textbf{Predeclared policy class:} $\Pi$ is fixed before looking at expensive privacy labels, or any adaptive expansion of $\Pi$ consumes additional multiplicity budget.
\item \textbf{Censoring discipline:} aborts, missing traces, failed lookups, and logging outages are either included in the SLA labels or counted as adverse outcomes under a declared rule.
\item \textbf{Replay determinism:} the transcript projection used for certification is the same projection later named by the public card; changing parser versions, padding semantics, or token roles creates a successor witness/interface rather than a wrapper refresh.
\end{enumerate}
We treat the estimator family as a deployment choice; the menu theorem only requires that each $\mathrm{UCB}$ be valid at its allocated level.

\section{Two-sample decoupling: expensive labels, cheap telemetry}
Many privacy labels require heavyweight analysis (e.g., reconstructing committee contacts or stop times).
A pragmatic trick is to certify a cheaper proxy statistic $Y$ on every run and only compute $X$ (the full label) on a subsample.
A two-sample protocol checks that $Y$ predicts $X$ with sufficient reliability to support certification.
(Operationally: publish both a label and a proxy audit bound in the receipt.)


\section{Deployment-facing validity checks}
Before this note is used as a public release root, the certified-menu card should answer five negative questions.
First, can an adversary change the observation surface without changing the witness/interface version?
Second, can an operator switch into a policy that was not in the declared raw pool?
Third, can missing or aborted attempts be dropped from the log?
Fourth, can off-policy support fail because a target route had zero or tiny logging probability?
Fifth, can a downstream wrapper reuse the old card after any menu-owned row changes?
A ``yes'' to any of these questions is not a packaging problem; it is a claim-validity problem.

Current deployed anonymity/privacy mechanisms are useful sanity checks because they expose exactly the operational seams that a mathematical menu card can accidentally abstract away.
Tor circuit padding, for example, is negotiated through padding machines and consensus parameters, and connection-level padding has consensus parameters that can disable or bound behavior.\cite{torpadding,torconnectionpadding}
Privacy Pass separates clients, origins, issuers, and attesters, so a token-flow deployment can change the witness/interface even when the high-level privacy story sounds unchanged.\cite{privacypass9576,privacypass9577}
For this paper, the safe rule is simple: if padding, retry, token, attester, issuer, origin, parser, or logging behavior changes a certified row, publish a successor card; if it only changes an outer ABOM/release wrapper while all menu-owned rows remain byte-identical, a wrapper refresh may cite the old card.

\section{Output format}
We recommend publishing:
(i) the certified menu policies (hashes/IDs),
(ii) their UCBs and thresholds,
(iii) the multiplicity accounting (how $\alpha$ was split or compressed), and
(iv) the exact witness/interface version against which the menu was certified.
ABOM/OINL (Anonymity~C) is a good carrier for the public claim surface, while Eval~1 can notarize the supporting attempt log and anti-grinding evidence.\cite{series:abomoinl,series:eval1}

\paragraph{A minimal public certified-menu card.}
Table~\ref{tab:minimal-certified-menu-card} is the smallest public packet later terse papers should need when the live question is simply \emph{which simultaneous menu contract was declared for this witness/interface version?}

\begin{table}[t]
\centering
\small
\begin{tabular}{@{}p{0.32\linewidth}p{0.59\linewidth}@{}}
\toprule
\textbf{Public row} & \textbf{Pinned value}\\
\midrule
Witness / interface version & committee-contact-count bucket on witness interface \texttt{cmenu-v1} \\
Raw policy pool & $M=8$ declared candidate policies $\{\pi_1,\ldots,\pi_8\}$ \\
Family-wise level and split & $\alpha=0.03$ split uniformly across $M$ policies and three SLA coordinates, so each coordinate gets $\alpha/(3M)=0.00125$ \\
SLA thresholds & privacy $\tau_P=0.80$; latency $\tau_T=120$ ms; bandwidth $\tau_B=4.0$ KB \\
Certified members & $\widehat\Pi=\{\pi_2,\pi_4,\pi_7\}$ under the declared raw-$M$ split \\
Representative bound row & for $\pi_4$: $\mathrm{UCB}_P=0.74$, $\mathrm{UCB}_T=118$ ms, $\mathrm{UCB}_B=3.9$ KB \\
Switchability claim & any later switch among $\pi_2,\pi_4,\pi_7$ stays inside the same certified menu contract \\
Downstream owner & Certified~B / Certified~C / Eval~1 / ABOM--OINL / Release~A for compression, drift, notarization, and lineage \\
\bottomrule
\end{tabular}
\caption{A minimal public certified-menu card. This is the non-decorative public answer later terse papers can quote directly when the live issue is the declared simultaneous menu contract rather than multiplicity compression, recertification, or release packaging.}
\label{tab:minimal-certified-menu-card}
\end{table}

This card is intentionally narrow: it pins the finite policy pool, the witness/interface version, the family-wise split, and the resulting switchable certified set, but nothing else.
Exact public/on-request packet wiring for one exported certified-menu card is intentionally delegated to the maintained worked example, and theorem-root-to-review routing across that already-fixed card is delegated to the maintained series-spine bridge.
Certified~B may later replace raw $M$ by an effective multiplicity $M_{\mathrm{eff}}$ for this same witness/interface version, while Certified~C, Eval~1, and Release~A own the drift, notarization, and lineage layers.\cite{series:workedexample,series:seriesspines,series:sigcomp,series:recert,series:eval1,series:releaseA}

\paragraph{One tiny raw-$M$ drill.}
With the card values above,
\[
\frac{\alpha}{3M}=\frac{0.03}{3\cdot 8}=0.00125.
\]
If the declared per-coordinate UCBs for $\pi_4$ are
\[
\mathrm{UCB}_P=0.74\le 0.80,\qquad \mathrm{UCB}_T=118\le 120,\qquad \mathrm{UCB}_B=3.9\le 4.0,
\]
then $\pi_4\in\widehat\Pi$ under the published raw-$M$ certificate.
A later terse paper can quote the card and this arithmetic without reopening off-policy estimation details or the release stack.

\paragraph{One tiny worked-route drill.}
Keep the same witness/interface version, raw policy pool, family-wise split, and certified-member row from Table~\ref{tab:minimal-certified-menu-card}.
Then a later terse paper may honestly say exactly one sentence: \emph{"Under witness/interface version \texttt{cmenu-v1}, the published certified-menu packet certifies the switchable set $\{\pi_2,\pi_4,\pi_7\}$ under the declared raw-$M$ split, so a switch to $\pi_4$ stays inside the same public menu contract."}
That sentence is honest because the maintained worked packet already fixes which public card carries the witness pin, split, and certified-member row, while theorem-root-to-review routing across that already-fixed card is maintained by the series-spine bridge rather than reconstructed locally here.\cite{series:workedexample,series:seriesspines}

\paragraph{One compact three-stop route.}
Suppose the maintained certified-menu card keeps witness/interface version \texttt{cmenu-v1}, the raw pool $M=8$, and the certified-member row $\{\pi_2,\pi_4,\pi_7\}$.
If that is the only live question, later terse papers should stop here and cite Certified~A.
If a separate signature-manifest card for the same witness/interface safely collapses a wider raw tuning pool from $M=12$ to $M_{\mathrm{eff}}=3$ on the declared weekday/weekend support, the honest second sentence belongs to Certified~B rather than to this note.
If a deployed update to $\pi_4$ later keeps the same ENF/TW/state identifiers and is accepted under a published recertification bound $\overline q=0.013\le \gamma=0.020$, the honest third sentence belongs to Certified~C rather than to this note.
That three-stop split is the stable owner map for certified routing: menu theorem/accountant root here, optional compression witness in Certified~B, and deployed change-control packet in Certified~C.\cite{series:sigcomp,series:recert}

\paragraph{A same-card / refreshed-wrapper cutover.}
The smallest public cutover after a certified-menu card exists is not merely semantic.
Suppose a later packet keeps the same witness/interface version \texttt{cmenu-v1}, the same raw policy pool $M=8$, the same family-wise split $\alpha/(3M)=0.00125$, the same SLA thresholds, and the same certified-member row $\widehat\Pi=\{\pi_2,\pi_4,\pi_7\}$, but only refreshes the outer ABOM or release-receipt wrapper.
Then the old certified-menu card remains the right public reference, and a later terse paper may keep citing it.
But if the successor changes any menu-owned row---the witness/interface pin, the raw policy pool, the $\alpha$ split, the SLA thresholds, the certified-member row, or the representative UCB row---then the old card no longer names the live simultaneous selection object, even if the prose still says ``same menu.''
That successor must therefore publish a fresh certified-menu card.
If the live change is merely that raw $M$ has been compressed to an imported effective multiplicity $M_{\mathrm{eff}}$, reopen Certified~B; if the live change is disagreement-triggered drift or recertification of a deployed menu, reopen Certified~C rather than stretching this baseline raw-$M$ note past its owner boundary.\cite{series:sigcomp,series:recert}

\paragraph{One tiny wrapper-refresh reuse drill.}
If a later release reissues the same \texttt{cmenu-v1} witness pin, the same eight-policy pool, the same $\alpha=0.03$ split, and the same certified-member row $\{\pi_2,\pi_4,\pi_7\}$, then a short paper may still say: \emph{"The maintained certified-menu card for \texttt{cmenu-v1} still licenses switches among $\pi_2,\pi_4$, and $\pi_7$; only the outer claim wrapper was refreshed."}
That sentence stops being honest as soon as the menu-owned rows drift, because at that point the later packet is not rewrapping the same simultaneous menu object---it is exporting a successor menu card.

\section{Takeaways}
\begin{itemize}[leftmargin=*]
\item \textbf{Selection is part of the threat model:} tuning across many configurations is a many-testing problem, so publish \emph{simultaneous} guarantees over a finite policy class rather than a single post-hoc ``best'' policy.
\item \textbf{Menus are switchable:} once a set is certified uniformly, operators can switch among policies in that set without re-running expensive experiments.
\item \textbf{Stop at the menu object:} this note does not own drift monitoring, evaluator notarization, or cross-release lineage; it exports the simultaneous menu object that those later papers consume.
\item \textbf{Do not confuse a valid UCB with a valid deployment claim:} off-policy support, censoring, logging propensities, and the adversary-facing transcript projection have to be pinned before the card is safe to reuse.
\item \textbf{The raw menu theorem already stands on its own:} publishing the raw policy count $M$, the $\alpha$ split, and the witness-version pin is enough for a complete referee-facing claim; Certified~B only improves tightness by compressing equivalent policies.
\item \textbf{The published claim still needs explicit multiplicity accounting:} whether the deployment uses the raw pool or an imported effective multiplicity, that accounting and the witness-version pin should be visible in the ABOM/OINL claim surface and any linked evaluation bundle.
\item \textbf{Reopen the right downstream owner:} if the live issue is the exact packet carrying one exported certified-menu card, reopen the maintained worked example; if it is theorem-root-to-review routing across that already-fixed card, reopen the maintained series-spine bridge rather than rebuilding either locally.
\end{itemize}

\begin{thebibliography}{99}\small

\bibitem{series:artifactpolicy}
Anonymous.
\newblock Artifact and Disclosure Policy for the One-True-Archive (Source-Only Public Release).
\newblock Series synthesis note (Synthesis~6), 2026. (This archive.)

\bibitem{series:oddsinflation}
Anonymous.
\newblock Odds Inflation, Privacy Loss, and Endpoint Anonymity Bounds.
\newblock Anonymity series Paper~A, 2026. (This archive.)

\bibitem{series:endpointbridge}
Anonymous.
\newblock Endpoint Metrics Bridge: Odds Inflation, PML/MaxL, and Identification Sizing.
\newblock Series synthesis note (Synthesis~5), 2026. (This archive.)

\bibitem{series:contractobjects}
Anonymous.
\newblock Contract Objects for Anonymous-DHT Privacy Budgets and Claims.
\newblock Series synthesis note (Synthesis~0), 2026. (This archive.)

\bibitem{series:abomoinl}
Anonymous.
\newblock ABOM/OINL: Audit BOM and Outcome Identification Noise Label for Anonymous DHT Deployments.
\newblock Anonymity series Paper~C, 2026. (This archive.)

\bibitem{series:state1}
Anonymous.
\newblock State-Dependent Anonymity: Selection and State Evolution as an Attack Surface.
\newblock AnonDHT State series Paper~1, 2026. (This archive.)

\bibitem{series:sigcomp}
Anonymous.
\newblock Routing-Signature Compression for Many-Testing-Safe Anonymous DHT Tuning.
\newblock Certified series Paper~B, 2026. (This archive.)

\bibitem{series:recert}
Anonymous.
\newblock Disagreement-Gated Recertification and Boundary Audits for Anonymous DHT Routers.
\newblock Certified series Paper~C, 2026. (This archive.)

\bibitem{series:eval1}
Anonymous.
\newblock Notarized Anonymity Certificates: Attempt Logs, Spending Discipline, and Verifier Rules.
\newblock Evaluation series Paper~1, 2026. (This archive.)

\bibitem{series:workedexample}
Anonymous.
\newblock Worked Example: Instantiating a Tiered Reader-Privacy Receipt.
\newblock Series synthesis note (Synthesis~17), 2026. (This archive.)

\bibitem{series:seriesspines}
Anonymous.
\newblock Series Spines from Mathematical Roots to Referee Handoffs.
\newblock Series synthesis note (Synthesis~55), 2026. (This archive.)

\bibitem{series:releaseA}
Anonymous.
\newblock Release Property Ledgers and Signed Receipts for Anonymous-DHT Privacy Claims.
\newblock Release and destination series Paper~A, 2026. (This archive.)


\bibitem{huang2021offpolicy}
A.~Huang, L.~Leqi, Z.~Lipton, and K.~Azizzadenesheli.
\newblock Off-policy risk assessment in contextual bandits.
\newblock \emph{NeurIPS}, 2021.

\bibitem{kuzborskij2021confident}
I.~Kuzborskij, C.~Vernade, A.~Gy\"orgy, and C.~Szepesv\'ari.
\newblock Confident off-policy evaluation and selection through self-normalized importance weighting.
\newblock \emph{AISTATS}, 2021.

\bibitem{torpadding}
The Tor Project.
\newblock Circuit-level padding specification.
\newblock \url{https://spec.torproject.org/padding-spec/circuit-level-padding.html}.

\bibitem{torconnectionpadding}
The Tor Project.
\newblock Connection-level padding specification.
\newblock \url{https://spec.torproject.org/padding-spec/connection-level-padding.html}.

\bibitem{privacypass9576}
IETF.
\newblock RFC 9576: The Privacy Pass Architecture.
\newblock \url{https://www.rfc-editor.org/info/rfc9576}, 2024.

\bibitem{privacypass9577}
IETF.
\newblock RFC 9577: The Privacy Pass HTTP Authentication Scheme.
\newblock \url{https://www.rfc-editor.org/rfc/rfc9577.html}, 2024.

\bibitem{holm1979}
S.~Holm.
\newblock A simple sequentially rejective multiple test procedure.
\newblock \emph{Scandinavian Journal of Statistics}, 6(2):65--70, 1979.

\end{thebibliography}

\end{document}
